% ============================================================
%  Booster T1 MJLab Port — Feature Gap Analysis & Goalkeeper Design
%  BEP Imitation Learning, 2026-05-24
% ============================================================
\documentclass[11pt,a4paper]{article}

\usepackage[margin=2.5cm]{geometry}
\usepackage{booktabs}
\usepackage{tabularx}
\usepackage{longtable}
\usepackage{xcolor}
\usepackage{listings}
\usepackage{hyperref}
\usepackage{amsmath}
\usepackage{parskip}
\usepackage{titlesec}
\usepackage{enumitem}
\usepackage{multirow}
\usepackage{array}
\usepackage{microtype}

% Colours
\definecolor{done}{RGB}{15,195,90}
\definecolor{partial}{RGB}{255,191,0}
\definecolor{codeblue}{RGB}{77,166,255}
\definecolor{codegrey}{RGB}{45,45,60}
\definecolor{lightgrey}{RGB}{240,240,245}
\definecolor{darkbg}{RGB}{26,26,46}

% Code listings style
\lstdefinestyle{python}{
  language=Python,
  basicstyle=\ttfamily\small,
  keywordstyle=\color{codeblue}\bfseries,
  commentstyle=\color{gray}\itshape,
  stringstyle=\color{done},
  numberstyle=\tiny\color{gray},
  numbers=left,
  numbersep=8pt,
  frame=single,
  framesep=4pt,
  backgroundcolor=\color{lightgrey},
  breaklines=true,
  breakatwhitespace=true,
  showstringspaces=false,
  tabsize=4,
  xleftmargin=1.5em,
  captionpos=b,
}
\lstset{style=python}

\hypersetup{
  colorlinks=true,
  linkcolor=black,
  urlcolor=blue,
  citecolor=black
}

% Section formatting
\titleformat{\section}{\large\bfseries}{}{0em}{}[\titlerule]
\titleformat{\subsection}{\normalsize\bfseries}{}{0em}{}

\newcommand{\done}{\textcolor{done}{\textbf{DONE}}}
\newcommand{\partial}{\textcolor{partial}{\textbf{PARTIAL}}}

% ============================================================
\begin{document}

\begin{center}
  {\LARGE\bfseries Booster T1 MJLab Port}\\[0.4em]
  {\large Feature Gap Analysis \& Goalkeeper Reward Design}\\[0.3em]
  {\normalsize BEP Imitation Learning --- 2026-05-24}
\end{center}

\vspace{0.5em}
\hrule
\vspace{1em}

\tableofcontents
\clearpage

% ============================================================
\section{Overview}

The original \textbf{InternRobotics Humanoid-Goalkeeper} pipeline was designed for
the Unitree G1 robot and runs inside NVIDIA Isaac Gym. This document describes the
gap analysis performed when porting the pipeline to the \textbf{Booster Robotics T1}
robot running inside MJLab (MuJoCo-based). Eight features were identified as missing
or incorrect in the port and subsequently implemented.

\subsection{Coordinate System Rotation}

The port applies a $+90^\circ$ yaw rotation throughout: the original has the ball
approach from $+X$; the port has the ball approach from $+Y$. All axis references
in this document use the \textbf{port convention} ($Y$ = approach axis, $X$ = lateral).

\subsection{Why Gaps Accumulated}

\begin{itemize}[nosep]
  \item $\approx$20 reactive debug sessions had been conducted without a full
        codebase comparison, so the same gaps were rediscovered repeatedly.
  \item The Isaac Gym \texttt{cfrc\_ext} contact-force tensor has no direct
        equivalent in MJLab, requiring explicit \texttt{ContactSensorCfg}
        declarations.
  \item Several curriculum mechanisms were embedded inside the Isaac Gym
        environment class (not in a config file) and were easy to miss.
\end{itemize}

The fix was a single exhaustive subagent pass that read \emph{all} upstream files
in one invocation, followed by parallel verification agents cross-checking each
claim, and then a single commit implementing all eight features.

% ============================================================
\section{Goalkeeper Reward Design}
\label{sec:goalkeeper}

This section explains how the goalkeeping task is structured so that training
behaviour can be interpreted correctly.

\subsection{Scenario}

\begin{itemize}[nosep]
  \item The ball is spawned at $y_\text{start} \in [3,5]$~m in front of the robot.
  \item The robot stands on the \textbf{goal line} ($y_\text{local} = 0$) and faces $+Y$.
  \item The ball is given an initial velocity aimed at a random target
        $x_\text{end} \in [-0.84,-0.2]$~m, $z_\text{end} \in [0.4,1.2]$~m
        (left-hand interception zone, full difficulty).
  \item Episode length: $3.0$~s at $50$~Hz = $150$ steps.
\end{itemize}

\subsection{Phase Structure}

The episode is divided into two phases by ball distance:

\begin{center}
\begin{tabular}{lll}
\toprule
\textbf{Phase} & \textbf{Condition} & \textbf{Primary reward} \\
\midrule
Pre-positioning & $y_\text{local} > 1.5$~m & \texttt{eereach} phase-1 (lateral alignment) \\
Interception    & $y_\text{local} \leq 1.5$~m & \texttt{eereach} phase-2 (sigmoid reach) \\
Post-save       & ball deflected or behind & \texttt{hand\_proximity\_strict}, \texttt{successland} \\
\bottomrule
\end{tabular}
\end{center}

\subsection{Stopball --- What It Actually Measures}
\label{subsec:stopball}

\texttt{stopball} is a \textbf{one-shot binary reward} that fires once per episode
when the ball's $Y$-velocity changes by $\geq 2$~m/s from its initial value:

\begin{lstlisting}[caption={stopball reward (MJLab port, \texttt{mdp/rewards.py})}]
def stopball(env, ball_name="ball", delta_vel_threshold=2.0):
    ball_y_vel   = ball.data.root_link_lin_vel_w[:, 1]
    ball_y_local = ball.data.root_link_pos_w[:, 1] - env_origins[:, 1]

    just_reset = env.episode_length_buf <= 1
    env._sb_flag[just_reset]    = False
    env._sb_init_vy[just_reset] = ball_y_vel[just_reset].clone()

    delta_vy = ball_y_vel - env._sb_init_vy   # current - initial
    in_front = ball_y_local > 0.0             # ball still in front of goal
    fired    = (delta_vy > delta_vel_threshold) & in_front & ~env._sb_flag

    env._sb_flag |= fired          # latch: prevents re-firing
    return fired.float()           # 0 or 1, maximum once per episode
\end{lstlisting}

\textbf{Key interpretation points:}

\begin{enumerate}[nosep]
  \item The threshold is a \emph{velocity change}, not an absolute velocity.
        A ball going from $-8$~m/s to $-5$~m/s satisfies $\Delta v_y = 3 > 2$
        and fires --- the ball is still moving at $5$~m/s.
  \item The task is \textbf{deflection}, not catching. A goalkeeper deflects;
        \texttt{hand\_proximity\_strict} incentivises keeping the hand on the
        ball after contact.
  \item Episode-sum of \texttt{stopball} = \textbf{1.0} means it fired exactly
        once --- this is the ceiling and indicates the robot reliably deflects
        the ball every episode.
  \item The \texttt{\_sb\_flag} latch prevents any episode-sum $> 1$.
        A logged value of \textbf{1.2} therefore suggests a double-fire bug
        (see Section~\ref{sec:stopball_bug}).
\end{enumerate}

\subsection{Complete Reward Table}
\label{subsec:rewards}

\begin{longtable}{p{3.5cm} r p{8.0cm}}
\toprule
\textbf{Term} & \textbf{Weight} & \textbf{Description} \\
\midrule
\endhead
\multicolumn{3}{l}{\textit{--- Task rewards ---}} \\[2pt]
\texttt{eereach}            & 10--20 & Sigmoid reach to pre-computed ball intercept point (\texttt{\_ball\_end\_target}). Phase-1: lateral pre-positioning when ball $> 1.5$~m. Phase-2: sigmoid $\times$ velocity amplification. \\[4pt]
\texttt{stopball}           & 100--200 & One-shot: $\Delta v_y > 2$~m/s while ball in front. Curriculum scales weight 100\,→\,150\,→\,200. \\[4pt]
\texttt{hand\_proximity\_strict} & 5--10 & Continuous: hand $< 0.15$~m from ball. Doubles to $2\times$ after \texttt{stopball} fires. Curriculum 5\,→\,7.5\,→\,10. \\[4pt]
\texttt{successland}        & 4 & Safe landing after jump save: +1 while airborne (\texttt{\_has\_in\_air}), +5 on landing, $-1$ one-foot landing. \\[4pt]
\midrule
\multicolumn{3}{l}{\textit{--- Stability / orientation ---}} \\[2pt]
\texttt{stayonline}         & $-2$ & Penalise $Y$ deviation from goal line $> 0.2$~m. \\[4pt]
\texttt{noretreat}          & $-2$ & Penalise retreating (negative body-frame forward velocity). \\[4pt]
\texttt{feetorientation}    & 3 & $\exp(-5 \sum (\text{foot gravity}_{xy})^2)$: feet flat. \\[4pt]
\texttt{postorientation}    & 3 & Upright posture reward active once ball is behind. \\[4pt]
\texttt{postangvel}         & 3 & Low angular velocity reward once ball is behind. \\[4pt]
\texttt{postlinvel}         & 1 & Low forward linear velocity reward once ball is behind. \\[4pt]
\texttt{postupperdofpos}    & 1 & Arms return to default pose once ball is behind: $\exp(-\sum\delta^2)$. \\[4pt]
\texttt{postwaistdofpos}    & 1 & Waist returns to default: $\exp(-3\sum\delta^2)$. \\[4pt]
\texttt{deviation\_waist}   & $-0.001$ & Always-active waist deviation penalty. \\[4pt]
\texttt{feet\_slippage}     & 3 & $\exp(-10 \sum v_\text{foot} \cdot \mathbf{1}_\text{contact})$: no sliding. \\[4pt]
\midrule
\multicolumn{3}{l}{\textit{--- Regularisation ---}} \\[2pt]
\texttt{action\_rate\_l2}   & $-0.1$ & Action rate (jerk) penalty. \\[4pt]
\texttt{dof\_acc}           & $-2.5\times10^{-7}$ & Joint acceleration L2. \\[4pt]
\texttt{torques}            & $-10^{-5}$ & Normalised torque L2 (divided by per-joint $k_p$). \\[4pt]
\texttt{dof\_vel}           & $-5\times10^{-4}$ & Joint velocity L2. \\[4pt]
\texttt{ang\_vel\_xy}       & $-0.1$ & Base angular velocity XY L2. \\[4pt]
\midrule
\multicolumn{3}{l}{\textit{--- Safety / limits ---}} \\[2pt]
\texttt{dof\_pos\_limits}   & $-3$ to $-9$ & Soft joint position limit violation (curriculum). \\[4pt]
\texttt{dof\_vel\_limits}   & $-2$ & Per-joint velocity limit violation. \\[4pt]
\texttt{torque\_limits}     & $-3$ to $-9$ & Per-joint torque limit violation (curriculum). \\[4pt]
\texttt{penalize\_sharpcontact} & $-100$ & Mean foot contact force $> 1000$~N. \\[4pt]
\texttt{penalize\_kneeheight}   & $-100$ & Shank body $< 0.12$~m above ground (kneeling). \\[4pt]
\texttt{penalize\_self\_collision} & $-50$ & Any Trunk-subtree self-collision detected. \\[4pt]
\midrule
\multicolumn{3}{l}{\textit{--- Motion tracking (from \texttt{make\_tracking\_env\_cfg}) ---}} \\[2pt]
\texttt{motion\_global\_root\_pos} & 4 & Global root position tracking. \\[4pt]
\texttt{motion\_global\_root\_ori} & 3 & Global root orientation tracking. \\[4pt]
\texttt{motion\_body\_pos}  & 4 & Body position tracking. \\[4pt]
\texttt{motion\_body\_ori}  & 2 & Body orientation tracking. \\[4pt]
\texttt{motion\_body\_lin\_vel} & 2 & Body linear velocity tracking. \\[4pt]
\texttt{motion\_body\_ang\_vel} & 2 & Body angular velocity tracking. \\[4pt]
\bottomrule
\caption{Complete reward configuration for Booster T1 goalkeeper training.
  Curriculum weights show initial\,→\,stage-1 (iter 600)\,→\,stage-2 (iter 1200).}
\end{longtable}

\subsection{Observation Space}

\begin{center}
\begin{tabular}{llll}
\toprule
\textbf{Observation} & \textbf{Dim} & \textbf{Actor} & \textbf{Critic} \\
\midrule
Joint positions (relative) & 21 & \checkmark & \checkmark \\
Joint velocities & 21 & \checkmark & \checkmark \\
Base angular velocity & 3 & \checkmark & \checkmark \\
Projected gravity & 3 & \checkmark & \checkmark \\
Actions (prev step) & 21 & \checkmark & \checkmark \\
Ball position (body frame, masked) & 3 & \checkmark & \checkmark \\
Ball velocity (body frame, masked) & 3 & \checkmark & \checkmark \\
Left hand position (body frame) & 3 & \checkmark & \checkmark \\
Right hand position (body frame) & 3 & \checkmark & \checkmark \\
Base linear velocity & 3 & $\times$ & \checkmark \\
\midrule
History length & & 10 & 10 \\
\bottomrule
\end{tabular}
\end{center}

Ball observations are masked by \texttt{\_compute\_ball\_visibility} during the
initial warmup (first $\approx 7$ steps) and randomly during flight (Feature 8).

% ============================================================
\section{8 Missing Features: Implementation Summary}
\label{sec:features}

\begin{longtable}{p{1.0cm} p{2.5cm} p{5.0cm} p{3.5cm} p{1.5cm}}
\toprule
\textbf{ID} & \textbf{Feature} & \textbf{What changed} & \textbf{Upstream ref} & \textbf{Status} \\
\midrule
\endhead
P1B & \texttt{successland} & Added \texttt{\_has\_in\_air} tracking; +5 landing bonus; $-1$ one-foot penalty; removed always-firing feet-down gate &
  \texttt{legged\_robot.py}:1445 & \done \\[6pt]

P1C & \texttt{penalize\-\_sharp\-contact} & Replaced trunk-height proxy with \texttt{feet\_contact} sensor; per-foot max then 2-body mean (matches upstream 2-body average) &
  \texttt{legged\_robot.py}:1475 & \done \\[6pt]

P2A & Ball difficulty curriculum & \texttt{\_BALL\_END\_RANGES\_EASY} starting range; \texttt{ball\_difficulty\_curriculum} class in \texttt{resets.py}; linear interpolation $0 \to 0.5 \to 1.0$ at iter 600/1200 &
  \texttt{legged\_robot.py}:366 & \done \\[6pt]

P2B & dof/torque limits curriculum & 3-stage weight scaling: $-3 \to -6 \to -9$ at iter 600/1200; mirrors upstream $\times 2, \times 3$ per curriculum step &
  \texttt{legged\_robot.py}:362 & \done \\[6pt]

P2C & \texttt{hand\_proximity\-\_strict} curriculum & 3-stage weight: $5 \to 7.5 \to 10$; mirrors upstream \texttt{success\_init} $\times (1 + 0.5 \times \text{update})$ &
  \texttt{legged\_robot.py}:1401 & \done \\[6pt]

P3A & eereach intercept target & Pre-computed \texttt{\_ball\_end\_target} at reset using $\texttt{catch\_prop} = (y_\text{start} - 0.1)/(y_\text{start} + 0.3)$; full phase-1 lateral alignment + phase-2 sigmoid &
  \texttt{legged\_robot.py}:797 & \done \\[6pt]

F7 & Catchstep warmup & \texttt{\_catchstep}$=50$ at reset; decremented per step; gates \texttt{ball\_pos\_b}/\texttt{ball\_vel\_b} visibility. \emph{Note: action-level override} (\texttt{joint\_pos\_target} override) \emph{not possible in MJLab without framework fork --- ball-observation masking used instead.} &
  \texttt{legged\_robot.py}:643 & \partial \\[6pt]

F8 & Ball visibility masking & \texttt{\_compute\_ball\_visibility}: (1) \texttt{initial\_vanish} ($\texttt{catchstep} < 43$), (2) \texttt{flying} (in-range, approaching, launched), (3) \texttt{random\_vanish} (per-env random step 0--30) &
  \texttt{legged\_robot.py}:968 & \done \\[6pt]
\bottomrule
\caption{Eight features ported from Humanoid-Goalkeeper to Booster T1 MJLab port.}
\end{longtable}

% ============================================================
\section{Stopball: Exact Upstream Implementation \& Analysis of 1.2}
\label{sec:stopball_bug}

\subsection{Exact original \texttt{\_reward\_stopball} (Isaac Gym, G1)}

Verified by subagent reading \texttt{legged\_robot.py} verbatim.
Lines 1407--1416:

\begin{lstlisting}[caption={Original stopball reward, legged\_robot.py:1407--1416}]
def _reward_stopball(self):
    changevel = (self.ball_states[:,7] - self.ball_vel > 2.0) \
              & ((self.ball_states[:, 0] - self.env_origins[:, 0]) > 0.0)
    stopped_ids = (changevel
        | ((self.ball_states[:, 0] - self.env_origins[:, 0]) < 0.0)
    ).nonzero(as_tuple=False).flatten()
    success_ids = ((self.stop_flag == 0) & changevel) \
                  .nonzero(as_tuple=False).flatten()
    rew_stop = 1.0 * (self.stop_flag == 0) * changevel
    self.success_flag[success_ids] = 1.0
    self.stop_flag[stopped_ids]    = 1.0
    return rew_stop
\end{lstlisting}

\textbf{Key facts from verbatim reading:}

\begin{itemize}[nosep]
  \item \texttt{ball\_states[:,7]} = ball linear velocity along $X$ (the approach axis
        in the original).  Index~7 in Isaac Gym's root state is $v_x$.
  \item \texttt{self.ball\_vel} is a \textbf{snapshot stored at episode reset} ---
        it stores only the initial $v_x$ component (line~713:
        \texttt{self.ball\_vel[ball\_ids] = ball\_vel[:,0]}). It is \emph{never
        updated mid-episode}.
  \item The threshold 2.0~m/s is \textbf{hardcoded everywhere} --- in
        \texttt{\_reward\_stopball}, \texttt{post\_physics\_step}, and
        \texttt{\_reward\_eereach}. It does not appear in \texttt{g1\_29\_config.py}.
  \item \texttt{stop\_flag} is set for \textbf{both} envs where the ball was
        deflected (\texttt{changevel}) \emph{and} envs where the ball went behind
        the robot (position $< 0$). This prevents re-firing even after a miss.
  \item \texttt{success\_flag} (used by \texttt{\_reward\_success}) is set
        \emph{only} for successful deflections --- it is a separate tensor.
  \item Double-fire is \textbf{impossible}: \texttt{rew\_stop} is gated by
        \texttt{(stop\_flag == 0)}, and \texttt{stop\_flag} is reset to~0
        only in \texttt{\_reset\_root\_states} (line~720).
\end{itemize}

\subsection{Exact original \texttt{\_reward\_success} (hand proximity)}

Lines 1402--1403:
\begin{lstlisting}[caption={Original hand proximity reward}]
def _reward_success(self):
    return (self.success_flag + 1.0) * (self.dist < self.cfg.rewards.strict_th)
\end{lstlisting}

Returns $1.0 \times (d < 0.15)$ before the stop; $2.0 \times (d < 0.15)$ after.
Fires \textbf{every step} the hand is within 0.15~m. There is no reward anywhere
for the ball having low absolute velocity after the deflection.

\subsection{Why the logged value is 1.2 (not a bug)}

The original reports rewards via (line~316):
\begin{lstlisting}[caption={Isaac Gym reward logging (legged\_robot.py:316)}]
self.extras["episode"]['rew_' + key] = torch.mean(
    self.episode_sums[key][env_ids]
    / torch.clip(self.episode_length_buf[env_ids], min=1)
    / self.dt
)
\end{lstlisting}

At startup, reward scales are multiplied by \texttt{dt} (line~1072:
\texttt{self.reward\_scales[key] *= self.dt}).
Additionally, stopball's scale is curriculum-scaled each iteration (lines~363--364):
\begin{lstlisting}
if "stopball" in self.reward_scales:
    self.reward_scales["stopball"] = self.stop_init * (1 + 0.5 * self.curriculumupdate)
\end{lstlisting}

So the logged value is:
\[
  \frac{\text{episode\_sum}}{\text{episode\_length} \times dt}
  = \frac{1.0 \times \text{scale} \times dt}{\text{episode\_length} \times dt}
  = \frac{\text{scale}}{\text{episode\_length}}
\]
A value of \textbf{1.2} reflects the current curriculum-scaled reward weight
divided by the episode length --- \textbf{not a double-fire}. The flag is
airtight.

For the MJLab port, MJLab logs the \emph{raw unweighted} function value summed
over the episode. A value of 1.2 there implies the function returns values
summing to 1.2 per episode on average. Since the raw function is always 0 or 1
and the flag prevents re-firing, the only remaining explanation is that
\texttt{episode\_length\_buf} behaviour causes the flag to be cleared twice
within one episode (see Section~\ref{subsec:off-by-one}).

\subsection{Port-specific flag bug: \texttt{stop\_flag} not set on miss}
\label{subsec:miss-flag}

The original sets \texttt{stop\_flag} for \textbf{both} deflections and misses
(ball behind goal). The port's \texttt{\_sb\_flag} is set \emph{only} when
\texttt{fired} is True (which requires \texttt{in\_front}):

\begin{lstlisting}[caption={Port: flag set only on deflection (potential gap)}]
fired = (delta_vy > delta_vel_threshold) & in_front & ~env._sb_flag
env._sb_flag |= fired   # NOT set when ball goes behind without deflection
\end{lstlisting}

If the ball passes through the goal without deflection, \texttt{\_sb\_flag}
stays False. In MuJoCo (no back wall), the ball continues behind the robot
indefinitely, so this is harmless in practice --- but the \texttt{in\_front}
gate in \texttt{fired} already prevents firing when \texttt{ball\_y\_local < 0}.

\subsection{Port-specific \texttt{episode\_length\_buf} off-by-one}
\label{subsec:off-by-one}

\begin{lstlisting}[caption={Port reset logic in stopball}]
just_reset = env.episode_length_buf <= 1   # True for steps 0 AND 1 if buf starts at 0
env._sb_flag[just_reset]    = False
env._sb_init_vy[just_reset] = ball_y_vel[just_reset].clone()
\end{lstlisting}

If MJLab's \texttt{episode\_length\_buf} starts at 0, this clears the flag on
\emph{two consecutive steps}. In practice the ball cannot be deflected by
$\geq 2$~m/s in the first two steps, so this is likely benign. However, a
cleaner fix stores the initial velocity directly in \texttt{\_reset\_ball}
when the launch velocity is known, rather than inferring it from physics steps:

\begin{lstlisting}[caption={Recommended fix: store initial vy at reset time}]
# In MultiMotionCommand._reset_ball (commands.py):
if not hasattr(self._env, "_sb_init_vy"):
    self._env._sb_init_vy = torch.zeros(self._env.num_envs, device=self.device)
    self._env._sb_flag    = torch.zeros(self._env.num_envs, dtype=torch.bool,
                                        device=self.device)
self._env._sb_init_vy[env_ids] = vy          # vy already computed above
self._env._sb_flag[env_ids]    = False
\end{lstlisting}

\subsection{Why the ball still has high velocity after stopball fires}

This is a \textbf{design limitation of the upstream reward}, not a port bug.
\texttt{stopball} fires on any $\Delta v_y > 2$~m/s change \emph{from the
initial launch velocity}. Concretely:

\begin{center}
\begin{tabular}{lllc}
\toprule
Initial $v_y$ & Final $v_y$ & $\Delta v_y$ & Fires? \\
\midrule
$-8$\,m/s & $-5.5$\,m/s & $+2.5$ & \textbf{Yes} --- ball still heading to goal \\
$-8$\,m/s & $0$\,m/s    & $+8$   & \textbf{Yes} --- ball stopped \\
$-8$\,m/s & $+2$\,m/s   & $+10$  & \textbf{Yes} --- ball bouncing back \\
\bottomrule
\end{tabular}
\end{center}

All three fire identically. Only the \texttt{in\_front} guard (ball still in
front of goal line) prevents fires after a miss. There is \textbf{no reward
anywhere} for the ball having low absolute velocity or bouncing back toward
$+Y$ (away from goal).

\subsection{Recommended fix: require ball to reverse direction}
\label{subsec:stopball_fix}

A goalkeeper that only grazes the ball (reducing $v_y$ from $-8$ to $-5$~m/s)
is not making a meaningful save --- the ball still enters the goal. A meaningful
save requires the ball to be deflected hard enough to leave the danger zone.
The correct condition is:

\begin{lstlisting}[caption={Improved stopball: require ball bounces back (positive $v_y$)}]
# Option A: ball must be going away from goal (positive vy)
bounced_back = ball_y_vel > 0.0
fired = bounced_back & in_front & ~env._sb_flag

# Option B: velocity change exceeds initial magnitude (direction reversed)
# Equivalent to ball_y_vel > 0 when init_vy < 0
fired = (delta_vy > torch.abs(env._sb_init_vy)) & in_front & ~env._sb_flag

# Option C: intermediate -- ball must slow to < 2 m/s (nearly stopped)
fired = (ball_y_vel > -2.0) & in_front & ~env._sb_flag
\end{lstlisting}

\textbf{Option A / B} is the most principled: the ball direction reversal
($v_y > 0$) is the unambiguous ``save'' signal for a goalkeeper blocking a shot
that approaches from $+Y$. It is equivalent to ``the ball is now moving away
from the goal''.

\textbf{Option C} is a softer version that allows the ball to slow to a
near-stop without full bounce-back, which is physically realistic for a heavy
ball absorbed by both hands.

\textbf{Consequence for training:} the current threshold is easy to satisfy
(any glancing contact that changes velocity by 2~m/s fires the reward), which
may explain why \texttt{stopball} saturates at 1 (or 1.2 in WandB) early in
training while the robot has not yet learned to make proper saves.

\subsubsection{Complementary: ball-arrest reward}

Even with an improved \texttt{stopball} trigger, there is no per-step incentive
to \emph{keep} the ball slow after it has been deflected. Adding:
\[
r_\text{arrest} = \mathbf{1}[\texttt{\_sb\_flag}] \cdot \exp\!\bigl(-k\,|v_y|\bigr)
\]
would reward the robot every step the ball is moving slowly after a save,
incentivising sustained hand pressure. Neither the upstream nor the current
port has this term.

% ============================================================
\section{Upstream vs Port: Reward Architecture Comparison}
\label{sec:arch}

\begin{center}
\begin{tabular}{p{4.5cm} p{4.5cm} p{4.5cm}}
\toprule
\textbf{Component} & \textbf{G1 upstream (Isaac Gym)} & \textbf{T1 port (MJLab)} \\
\midrule
Physics backend & Isaac Gym (GPU parallel) & MuJoCo (MJLab) \\
Training algorithm & HIM-PPO + 6$\times$ AMP & Standard PPO (mjlab) \\
Motion prior & AMP discriminators (40\% reward) & Explicit motion tracking \\
Contact forces & \texttt{cfrc\_ext} per body & \texttt{ContactSensorCfg} \\
Ball approach axis & $+X$ & $+Y$ (90° rotation) \\
\midrule
eereach & Phase 1 + Phase 2 + vel\_sigma & Phase 1 + Phase 2 + vel\_sigma \checkmark \\
stopball & $\Delta v_x > 2$~m/s, weight 100 & $\Delta v_y > 2$~m/s, weight 100 \checkmark \\
hand\_proximity & weight 5, scales with curriculum & weight 5--10 (curriculum) \checkmark \\
successland & \_has\_in\_air; jump regions only & \_has\_in\_air; all envs \checkmark \\
sharpcontact & 2 ankle-roll bodies, force & 4 foot geoms (max per foot, mean) \checkmark \\
Ball curriculum & \texttt{command\_ranges} per env & \texttt{\_ball\_difficulty} float (lerp) \checkmark \\
Catchstep warmup & Action override (joint targets) & Ball-obs masking only (partial) \\
Ball visibility & 3-phase (G1 reference) & 3-phase \checkmark \\
\bottomrule
\end{tabular}
\end{center}

% ============================================================
\section{Methodology: How to Avoid Re-Discovering the Same Gaps}
\label{sec:method}

The following workflow was developed after 20+ reactive sessions failed to
achieve stable coverage of all upstream features:

\begin{enumerate}
  \item \textbf{Exhaustive subagent}: Spawn one agent with explicit instruction to
        read \emph{every file} in the upstream codebase in one pass. Forbid
        summarisation --- require full code quotes with line numbers.
  \item \textbf{Parallel verification}: For each identified gap, spawn a separate
        agent to cross-check the claim against the upstream source (4 agents in
        parallel). Only implement after verification.
  \item \textbf{Atomic commit}: Implement all verified gaps in one commit together
        with the \texttt{DIVERGENCE\_FROM\_UPSTREAM.md} documentation entry.
  \item \textbf{Memory-indexed reference}: Store the critical file map in
        \texttt{.claude/projects/.../memory/} so future sessions start with
        known locations, not blank exploration.
\end{enumerate}

\noindent Example subagent prompt pattern:

\begin{verbatim}
Read ALL files under <path>. Do not summarise any function.
Quote every relevant function verbatim with line numbers.
Answer: [specific question list].
\end{verbatim}

% ============================================================
\end{document}
